% ============================================================
\chapter{Vari\'et\'es Riemanniennes}
\label{ch:varietes}
% ============================================================

En 1854, Bernhard Riemann, alors \^ag\'e de seulement 27 ans, prononce sa
le\c{c}on d'habilitation \`a G\"ottingen devant un jury qui comprend Gauss
lui-m\^eme. Le titre~: \emph{\"Uber die Hypothesen, welche der Geometrie zu
Grunde liegen} --- \og Sur les hypoth\`eses qui servent de fondement \`a la
g\'eom\'etrie\fg. En une heure, Riemann bouleverse les fondements de la
g\'eom\'etrie~: il propose de remplacer l'espace euclidien rigide par des
\emph{espaces courbes}, o\`u la notion de distance varie d'un point \`a
l'autre. Gauss, dit-on, en fut profond\'ement \'emu.

L'id\'ee centrale de Riemann est \'el\'egante~: au lieu de d\'efinir une distance
globale, on se donne en chaque point un \emph{produit scalaire} sur
l'espace tangent. \`A partir de cette donn\'ee infinit\'esimale, on peut
reconstruire des distances, des angles, des aires, et m\^eme la courbure de
l'espace. C'est cette structure --- une vari\'et\'e munie d'un tel champ de
produits scalaires --- que nous appelons une \emph{vari\'et\'e riemannienne}.

% ------------------------------------------------------------
\section{M\'etriques riemanniennes}
% ------------------------------------------------------------

Formalisons l'intuition de Riemann. Sur une vari\'et\'e diff\'erentielle $M$,
chaque point $p$ poss\`ede un espace tangent $T_pM$, qui est un espace
vectoriel de dimension $n$. Une m\'etrique riemannienne consiste \`a choisir
un produit scalaire sur chacun de ces espaces tangents, de mani\`ere lisse.

\begin{definition}[M\'etrique riemannienne]
Soit $M$ une vari\'et\'e diff\'erentielle lisse de dimension $n$. Une \emph{m\'etrique
riemannienne} sur $M$ est la donn\'ee, pour chaque point $p \in M$, d'un produit
scalaire $g_p : T_pM \times T_pM \to \R$, variant de mani\`ere lisse~: pour tous
champs de vecteurs lisses $X, Y \in \mathfrak{X}(M)$, l'application
$p \mapsto g_p(X_p, Y_p)$ est $C^\infty$.

Le couple $(M, g)$ est appel\'e \emph{vari\'et\'e riemannienne}.
\end{definition}

En coordonn\'ees locales $(x^1, \ldots, x^n)$, la m\'etrique s'\'ecrit~:
\[
    g = g_{ij} \, dx^i \otimes dx^j, \qquad
    g_{ij} = g\!\left(\frac{\partial}{\partial x^i}, \frac{\partial}{\partial x^j}\right).
\]
La matrice $(g_{ij}(p))$ est sym\'etrique d\'efinie positive en tout point $p$.

\begin{theoreme}[Existence de m\'etriques riemanniennes]
Toute vari\'et\'e diff\'erentielle paracompacte admet une m\'etrique riemannienne.
\end{theoreme}

\begin{proof}
Soit $\{(U_\alpha, \varphi_\alpha)\}$ un atlas de $M$ et $\{\rho_\alpha\}$ une
partition de l'unit\'e subordonn\'ee au recouvrement $\{U_\alpha\}$. Sur chaque
ouvert de carte $U_\alpha$, on d\'efinit la m\'etrique euclidienne pullback~:
$g_\alpha = \varphi_\alpha^*(\delta_{ij}\, dx^i \otimes dx^j)$.
Alors $g = \sum_\alpha \rho_\alpha \, g_\alpha$ est une m\'etrique riemannienne
sur $M$, car une combinaison convexe de produits scalaires est un produit scalaire.
\end{proof}

\begin{remarque}
L'unicit\'e n'a pas lieu~: l'espace des m\'etriques riemanniennes sur $M$ est un
c\^one convexe ouvert de dimension infinie dans $\Gamma(S^2 T^*M)$.
\end{remarque}

\section{Premiers exemples fondamentaux}

\begin{exemple}[Espace euclidien]
L'espace $\R^n$ muni de la m\'etrique standard $g = \delta_{ij}\, dx^i \otimes dx^j$
est la vari\'et\'e riemannienne la plus simple. On a $g_{ij} = \delta_{ij}$ (symbole
de Kronecker).
\end{exemple}

\begin{exemple}[Sph\`ere $S^n$]
\label{ex:sphere}
La sph\`ere unit\'e $S^n = \{x \in \R^{n+1} : \norm{x} = 1\}$ est munie de la
m\'etrique induite par l'inclusion $\iota : S^n \hookrightarrow \R^{n+1}$~:
$g_{S^n} = \iota^* g_{\R^{n+1}}$.

En coordonn\'ees sph\'eriques sur $S^2$, avec $(\theta, \phi) \in (0,\pi) \times (0,2\pi)$~:
\[
    g_{S^2} = d\theta^2 + \sin^2\!\theta \, d\phi^2.
\]
La matrice de la m\'etrique est~:
$\displaystyle (g_{ij}) = \begin{pmatrix} 1 & 0 \\ 0 & \sin^2\!\theta \end{pmatrix}$,
de d\'eterminant $\det(g_{ij}) = \sin^2\!\theta$.
\end{exemple}

\begin{exemple}[Espace hyperbolique $\mathbb{H}^n$]
\label{ex:hyperbolique}
Le \emph{mod\`ele du demi-espace sup\'erieur}~: $\mathbb{H}^n = \{(x^1, \ldots, x^n)
\in \R^n : x^n > 0\}$ muni de~:
\[
    g_{\mathbb{H}^n} = \frac{(dx^1)^2 + \cdots + (dx^n)^2}{(x^n)^2}.
\]
Pour $n = 2$~: $g = \frac{dx^2 + dy^2}{y^2}$ avec $(x, y) \in \R \times \R_{>0}$.

Le \emph{mod\`ele du disque de Poincar\'e}~: $\mathbb{D}^n = \{x \in \R^n : \norm{x} < 1\}$
muni de~:
\[
    g_{\mathbb{D}^n} = \frac{4\left((dx^1)^2 + \cdots + (dx^n)^2\right)}{(1 - \norm{x}^2)^2}.
\]
Ces deux mod\`eles sont isom\'etriques.
\end{exemple}

\begin{exemple}[Tore plat]
\label{ex:tore}
Le tore $\mathbb{T}^n = \R^n / \Z^n$ h\'erite de la m\'etrique euclidienne standard
par passage au quotient. C'est une vari\'et\'e compacte plate (courbure nulle).
Plus g\'en\'eralement, pour un r\'eseau $\Lambda \subset \R^n$, le tore $\R^n/\Lambda$
est une vari\'et\'e riemannienne plate.
\end{exemple}

\begin{exemple}[Groupes de Lie et m\'etriques bi-invariantes]
\label{ex:lie}
Soit $G$ un groupe de Lie et $\mathfrak{g} = T_eG$ son alg\`ebre de Lie. Un
produit scalaire $\inner{\cdot}{\cdot}$ sur $\mathfrak{g}$ d\'efinit une m\'etrique
invariante \`a gauche par~:
\[
    g_a(u, v) = \inner{(dL_{a^{-1}})_a\, u}{(dL_{a^{-1}})_a\, v},
    \quad u, v \in T_aG,
\]
o\`u $L_a : G \to G$ est la translation \`a gauche.

La m\'etrique est \emph{bi-invariante} si elle est aussi invariante \`a droite, ce
qui \'equivaut \`a~:
\[
    \inner{\operatorname{ad}_X Y}{Z} + \inner{Y}{\operatorname{ad}_X Z} = 0,
    \quad \forall\, X, Y, Z \in \mathfrak{g}.
\]
C'est le cas de tout groupe de Lie compact (la forme de Killing chang\'ee de signe
convient pour les groupes semi-simples compacts).
\end{exemple}

\section{Longueur, distance et volume}

\begin{definition}[Longueur d'une courbe]
Soit $\gamma : [a,b] \to M$ une courbe lisse par morceaux. Sa \emph{longueur} est~:
\[
    L(\gamma) = \int_a^b \norm{\dot{\gamma}(t)} \, dt
    = \int_a^b \sqrt{g_{\gamma(t)}(\dot{\gamma}(t), \dot{\gamma}(t))} \, dt.
\]
\end{definition}

\begin{definition}[Distance riemannienne]
La \emph{distance riemannienne} entre deux points $p, q \in M$ est~:
\[
    d(p,q) = \inf \{ L(\gamma) : \gamma \text{ courbe lisse par morceaux de } p \text{ \`a } q \}.
\]
\end{definition}

\begin{theoreme}[Structure d'espace m\'etrique]
La distance $d$ d\'efinit sur $M$ une structure d'espace m\'etrique dont la topologie
co\"incide avec la topologie de vari\'et\'e.
\end{theoreme}

\begin{proof}
La s\'eparation $d(p,q) > 0$ pour $p \neq q$ d\'ecoule de l'utilisation de
coordonn\'ees locales~: dans une carte, la m\'etrique riemannienne est \'equivalente
\`a la m\'etrique euclidienne, donc toute courbe joignant $p \neq q$ a une longueur
strictement positive. L'in\'egalit\'e triangulaire est imm\'ediate par concat\'enation
de courbes. La compatibilit\'e topologique se d\'emontre en comparant les boules
m\'etriques avec les boules de coordonn\'ees.
\end{proof}

\begin{definition}[\'El\'ement de volume]
La \emph{forme volume riemannienne} est la $n$-forme d\'efinie en coordonn\'ees
locales orient\'ees par~:
\[
    dV_g = \sqrt{\det(g_{ij})} \, dx^1 \wedge \cdots \wedge dx^n.
\]
Le volume d'un domaine $\Omega \subset M$ est $\operatorname{Vol}(\Omega) = \int_\Omega dV_g$.
\end{definition}

\begin{formules}[title=\textbf{Formules de volume des espaces mod\`eles}]
\begin{align*}
    \operatorname{Vol}(S^n(r)) &= \frac{2\pi^{(n+1)/2}}{\Gamma\!\left(\frac{n+1}{2}\right)} r^n, \\
    \operatorname{Vol}(B_{\R^n}(0,r)) &= \frac{\pi^{n/2}}{\Gamma\!\left(\frac{n}{2}+1\right)} r^n, \\
    \operatorname{Vol}(B_{\mathbb{H}^n}(p,r)) &= \omega_{n-1} \int_0^r \sinh^{n-1}(t)\, dt.
\end{align*}
\end{formules}

\section{Isom\'etries et applications conformes}

\begin{definition}[Isom\'etrie]
Soient $(M, g)$ et $(N, h)$ deux vari\'et\'es riemanniennes. Un diff\'eomorphisme
$\varphi : M \to N$ est une \emph{isom\'etrie} si $\varphi^* h = g$, c'est-\`a-dire~:
\[
    h_{\varphi(p)}(d\varphi_p(u), d\varphi_p(v)) = g_p(u, v),
    \quad \forall\, p \in M,\; u, v \in T_pM.
\]
\end{definition}

\begin{theoreme}[Groupe d'isom\'etries -- Myers--Steenrod]
Le groupe des isom\'etries $\operatorname{Isom}(M, g)$ d'une vari\'et\'e riemannienne
est un groupe de Lie, de dimension au plus $\frac{n(n+1)}{2}$.
\end{theoreme}

\begin{remarque}
L'\'egalit\'e $\dim \operatorname{Isom}(M,g) = \frac{n(n+1)}{2}$ est atteinte si et
seulement si $(M, g)$ est un espace \`a courbure constante.
\end{remarque}

\begin{exemple}[Isom\'etries de $S^n$]
$\operatorname{Isom}(S^n) = O(n+1)$, de dimension $\frac{n(n+1)}{2}$.
\end{exemple}

\begin{exemple}[Isom\'etries de $\mathbb{H}^n$]
Dans le mod\`ele du demi-espace, pour $n=2$, les isom\'etries pr\'eservant l'orientation
sont les transformations de M\"obius $z \mapsto \frac{az+b}{cz+d}$ avec
$a, b, c, d \in \R$, $ad - bc = 1$, soit $\operatorname{Isom}^+(\mathbb{H}^2)
\cong PSL(2, \R)$.
\end{exemple}

\begin{definition}[Application conforme]
Un diff\'eomorphisme $\varphi : (M,g) \to (N,h)$ est \emph{conforme} s'il existe
$\lambda \in C^\infty(M)$, $\lambda > 0$, tel que $\varphi^* h = \lambda^2 g$.
\end{definition}

\section{Produits et rev\^etements riemanniens}

\begin{definition}[M\'etrique produit]
Soient $(M_1, g_1)$ et $(M_2, g_2)$ deux vari\'et\'es riemanniennes. Le
\emph{produit riemannien} est $(M_1 \times M_2, g_1 \oplus g_2)$, o\`u~:
\[
    (g_1 \oplus g_2)_{(p_1, p_2)}((u_1, u_2), (v_1, v_2))
    = (g_1)_{p_1}(u_1, v_1) + (g_2)_{p_2}(u_2, v_2).
\]
\end{definition}

\begin{exemple}
Le cylindre $S^1 \times \R$ et le tore $S^1 \times S^1$ sont des produits riemanniens.
Le tore plat $\mathbb{T}^2 = S^1(r_1) \times S^1(r_2)$ poss\`ede une courbure nulle.
\end{exemple}

\begin{definition}[Produit tordu]
Soit $(B, g_B)$ une base et $(F, g_F)$ une fibre. Pour $f \in C^\infty(B)$, $f > 0$,
le \emph{produit tordu} est $B \times_f F$ muni de~:
\[
    g = g_B + f^2 \, g_F.
\]
\end{definition}

\begin{exemple}
L'espace $\R^n \setminus \{0\}$ en coordonn\'ees polaires s'\'ecrit comme un produit
tordu $\R_{>0} \times_r S^{n-1}$ avec $g = dr^2 + r^2 g_{S^{n-1}}$.
\end{exemple}

\begin{proposition}[Rev\^etement riemannien]
Si $\pi : \widetilde{M} \to M$ est un rev\^etement et $(M, g)$ une vari\'et\'e
riemannienne, alors il existe une unique m\'etrique $\tilde{g}$ sur $\widetilde{M}$
telle que $\pi$ soit une isom\'etrie locale~: $\pi^* g = \tilde{g}$.
\end{proposition}

\section{Compl\'etude}

\begin{definition}[Compl\'etude m\'etrique]
$(M, g)$ est \emph{compl\`ete} si l'espace m\'etrique $(M, d)$ est complet au sens
de Cauchy, c'est-\`a-dire si toute suite de Cauchy converge.
\end{definition}

\begin{theoreme}[Hopf--Rinow]
\label{thm:hopf_rinow}
Soit $(M, g)$ une vari\'et\'e riemannienne connexe. Les conditions suivantes sont
\'equivalentes~:
\begin{enumerate}
    \item $(M, d)$ est complet comme espace m\'etrique.
    \item Les sous-ensembles ferm\'es born\'es sont compacts.
    \item Il existe un point $p \in M$ tel que $\exp_p$ est d\'efini sur tout $T_pM$.
    \item Pour tout $p \in M$, $\exp_p$ est d\'efini sur tout $T_pM$.
\end{enumerate}
De plus, si l'une de ces conditions est v\'erifi\'ee, alors deux points quelconques
de $M$ sont joints par une g\'eod\'esique minimisante.
\end{theoreme}

\begin{attention}[title=\textbf{Compl\'etude et compacit\'e}]
Compl\'etude n'implique \emph{pas} compacit\'e~: $\R^n$ et $\mathbb{H}^n$ sont complets
mais non compacts. En revanche, toute vari\'et\'e riemannienne compacte est compl\`ete.
\end{attention}

\section{Calculs en coordonn\'ees}

\begin{exemple}[M\'etrique en coordonn\'ees polaires sur $\R^2$]
Le changement de coordonn\'ees $x = r\cos\theta$, $y = r\sin\theta$ donne~:
\[
    g = dx^2 + dy^2 = dr^2 + r^2\, d\theta^2.
\]
D'o\`u $(g_{ij}) = \begin{pmatrix} 1 & 0 \\ 0 & r^2 \end{pmatrix}$ et
$\sqrt{\det g} = r$.
\end{exemple}

\begin{exemple}[M\'etrique de Fubini--Study sur $\C P^1$]
On identifie $\C P^1 \cong S^2$. En coordonn\'ees inhomog\`enes $z = x + iy$~:
\[
    g_{FS} = \frac{4\,(dx^2 + dy^2)}{(1 + x^2 + y^2)^2}
    = \frac{4\, |dz|^2}{(1 + |z|^2)^2}.
\]
C'est la m\'etrique \`a courbure constante $K = 1$ sur $S^2$.
\end{exemple}

\section{Fonctorielles de la m\'etrique}

\begin{proposition}[Pullback de m\'etrique]
Si $f : M \to N$ est une immersion et $(N, h)$ est riemannienne, alors
$g = f^* h$ est une m\'etrique riemannienne sur $M$ (appel\'ee \emph{m\'etrique induite}).
\end{proposition}

\begin{proposition}[Op\'erateurs musicaux]
La m\'etrique induit des isomorphismes canoniques~:
\begin{align*}
    \flat : TM &\to T^*M, & v &\mapsto g(v, \cdot), \\
    \sharp : T^*M &\to TM, & \omega &\mapsto \text{l'unique } v \text{ tel que }
    g(v, \cdot) = \omega.
\end{align*}
En coordonn\'ees~: $(v^\flat)_i = g_{ij} v^j$ et $(\omega^\sharp)^i = g^{ij} \omega_j$,
o\`u $(g^{ij})$ est la matrice inverse de $(g_{ij})$.
\end{proposition}

\section{Exercices}

\begin{exercice}
Montrer que la m\'etrique hyperbolique $g = \frac{dx^2 + dy^2}{y^2}$ sur
$\mathbb{H}^2$ donne la distance~:
\[
    d((x_1, y_1), (x_2, y_2))
    = \operatorname{arcosh}\!\left(1 + \frac{(x_1-x_2)^2 + (y_1-y_2)^2}{2 y_1 y_2}\right).
\]
\end{exercice}

\begin{exercice}
Calculer le volume de la sph\`ere $S^2(r)$ de rayon $r$ en int\'egrant la forme volume
en coordonn\'ees sph\'eriques. V\'erifier que $\operatorname{Vol}(S^2(r)) = 4\pi r^2$.
\end{exercice}

\begin{exercice}
Soit $G = SO(3)$ muni de la m\'etrique bi-invariante $g(X, Y) = -\frac{1}{2}\tr(XY)$
pour $X, Y \in \mathfrak{so}(3)$. Montrer que cette m\'etrique fait de $SO(3)$ un
espace \`a courbure sectionnelle constante $K = 1/4$.
\end{exercice}

\begin{exercice}
Montrer que le produit tordu $M = \R \times_{\cosh t} \R$ muni de
$g = dt^2 + \cosh^2(t)\, ds^2$ est isom\'etrique \`a un ouvert de $\mathbb{H}^2$.
\end{exercice}

\begin{exercice}
Soit $(M, g)$ une vari\'et\'e riemannienne et $\lambda \in C^\infty(M)$, $\lambda > 0$.
On pose $\tilde{g} = \lambda^2 g$ (changement conforme). Calculer la relation entre
les formes volume $dV_{\tilde{g}}$ et $dV_g$.
\end{exercice}
